\fancyhead[C]{\textit{Solutions}}

\section{Solutions to Section 1}
	\begin{table}[!h]
		\centering
		\begin{tabular}{c@{\hspace{2.5em}}cp{0.75\textwidth}}
			true & false & \\
			\specialrule{\lightrulewidth}{0.25em}{1.25em}
			\bubble & \dbubble & For any scalar \(c \in \mathbb{R}\), \(\|c\vec{v}\| = c\|\vec{v}\|\). \\[0.75em]
			\dbubble & \bubble & If \(\vec{u} \cdot \vec{v} = 0\), then \(\|\vec{v} - \vec{u}\| = \|\vec{v} + \vec{u}\|.\)\\[0.75em]
			\bubble & \dbubble & If \(\vec{y} \in \mathbb{R}^n\) is orthogonal to some \(\vec{z} \in W\), 
						      where \(W\) is a subspace of \(\mathbb{R}^n\), then \(\vec{y} \in W^\perp\). \\[0.75em]
			\bubble & \dbubble & If \(A \in \mathbb{R}^{10 \times 12}\) and \(\rank A = 6\), then \(\dim (\nullspace A^T) = 6\).\\[0.75em]
			\bottomrule
		\end{tabular}
	\end{table}

\subsection{Justifications}
\paragraph{Justification to \textit{Q1}} It should be \(\lvert c\rvert \|\vec{v}\|\).

\paragraph{Justification to \textit{Q2}} Calculate \(\|\vec{v} - \vec{u}\|^2\) and \(\|\vec{v} + \vec{u}\|^2\). Since the norm always defines a positive value, 
							we know that if the squares are equal, then so are the square roots.

\begin{align*}
	\|\vec{v} - \vec{u}\|^2 &= (\vec{v} - \vec{u})^T(\vec{v} - \vec{u}) \\
				           &= \vec{v}^T\vec{v} - \vec{v}^T\vec{u} - \vec{u}^T\vec{v} + \vec{u}^T\vec{u} \\
                                           &= \|\vec{v}\|^2 + \|\vec{u}\|^2
\end{align*}

\begin{align*}
	\|\vec{v} + \vec{u}\|^2 &= (\vec{v} + \vec{u})^T(\vec{v} + \vec{u}) \\
				           &= \vec{v}^T\vec{v} + \vec{v}^T\vec{u} + \vec{u}^T\vec{v} + \vec{u}^T\vec{u} \\
                                           &= \|\vec{v}\|^2 + \|\vec{u}\|^2
\end{align*}

Since \(\|\vec{v} - \vec{u}\|^2 = \|\vec{v} + \vec{u}\|^2\), \(\|\vec{v} - \vec{u}\| = \|\vec{v} + \vec{u}\|\).

\paragraph{Justification to \textit{Q3}} The vector \(\vec{y}\) has to be orthogonal to \emph{every} vector in \(W\), not just \emph{some}.

\paragraph{Justification to \textit{Q4}} \(\dim (\nullspace A^T) = 4\).


\section{Solutions to Section 2}
	\begin{table}[!h]
		\centering
		\begin{tabular}{c@{\hspace{2.5em}}cp{0.75\textwidth}}
			possible & impossible & \\
			\specialrule{\lightrulewidth}{0.25em}{1.25em}
			\dbubble & \bubble & A vector whose projection onto a subspace \(W\) is itself. \\[0.75em]
			\dbubble & \bubble & A non-zero vector whose projection into a subspace is the zero vector.\\[0.75em]
			\bubble & \dbubble & 2 orthogonal planes in \(\mathbb{R}^3\). \\[0.75em]
			\dbubble & \bubble & Two vectors that form a linearly dependent set but also an orthogonal set. \\[0.75em]
			\bubble & \dbubble & \(A\) is orthogonal matrix and \(B\) is an orthogonal matrix as well, but \(AB\) isn't. \\[0.75em]
			\bottomrule
		\end{tabular}
	\end{table}
	
	\subsection{Justifications}
	\paragraph{Justification to Q1} For any vector \(\vec{y} \in W\), where \(W\) is a subspace, \(\operatorname{proj}_W \vec{y} = \vec{y}\).
	\paragraph{Justification to Q2} Any orthogonal vector to \(W\) will have a projection equal to the zero vector.
	\paragraph{Justification to Q3} The dimension of orthogonal complements have to add up to the dimension of the vector space, which is 3 in this case.
	\paragraph{Justification to Q4} \((AB)(AB)^T = A(BB^T)A^T = AA^T = I.\) So \(AB\) is orthogonal as well.

	\subsection{Solution to Question 3}
	Note that the matrix \(A\) is orthogonal. Hence, \(\|A\vec{v}\| = \|\vec{v}\|.\) We see that \(\|\vec{v}\| = \sqrt{10} = \|A\vec{v}\|\).

\section{Solutions to Section 3}
	\begin{table}[!h]
		\centering
		\begin{tabular}{c@{\hspace{2.5em}}cp{0.75\textwidth}}
			possible & impossible & \\
			\specialrule{\lightrulewidth}{0.25em}{1.25em}
			\bubble & \dbubble & The projection of a vector onto a subspace depends on the basis of the subspace. \\[0.75em]
			\bubble & \dbubble & A nonzero vector \(\vec{v}\) with \(\vec{v} \in W\) and \(\vec{v} \in W^\perp\). \\[0.75em]
			\bubble & \dbubble & A vector \(\vec{y}\) such that \(\vec{y} - \operatorname{proj}_W \vec{y} \notin W^\perp\). \\[0.75em]
			\bottomrule
		\end{tabular}
	\end{table}

	\subsection{Justifications}
	\paragraph{Justification to Q1} The projection of a vector onto a subspace is unique and does not depend on the basis of the subspace but rather the 
						      subspace itself.
	\paragraph{Justification to Q2} If \(\vec{v} \in W\) and \(\vec{v} \in W^\perp\), then \(\vec{v} \cdot \vec{v} = 0\). The only vector with this property is the
						      zero vector.
	\paragraph{Justification to Q3} Refer back to how orthogonal decomposition works.


	\subsection{Solution to Question 5}
	\begin{enumerate}
		\item \textbf{3}. A picture.
		\begin{figure}[!h]
			\centering
			\begin{tikzpicture}
				\begin{axis}[height = 5cm, width=8cm, axis lines = middle, xlabel = {\(x_1\)}, ylabel = {\(x_2\)}, xtick distance = 1, ytick distance = 1, xmin=0, xmax=5, ymin=0, ymax=5, grid=major]
					\plot[blue, ->]coordinates {(0, 0) (2, 3)} node[right, fill=white, scale=0.75] {\(\begin{pmatrix}2 \\ 3\end{pmatrix}\)};
					\plot[<->, semithick, purple, domain=0:5] ({x}, {0});
					\draw[semithick, ->] (axis cs: 2, 0) -- node[right, fill=white, scale=0.75] {\(\biggl\|\begin{pmatrix}2 \\ 3\end{pmatrix} - \operatorname{proj}_{\vec{e}_1} \begin{pmatrix}2 \\ 3\end{pmatrix}\biggr\| = 3\)} (axis cs: 2, 3) ;
				\end{axis}
			\end{tikzpicture}
		\end{figure}
		\item \(\mathbf{\sqrt{10}}\). Another picture.
		\begin{figure}[!t]
			\centering
			\begin{tikzpicture}
				\begin{axis}[height = 8cm, width=8cm, axis lines = middle, xlabel = {\(x_1\)}, ylabel = {\(x_2\)}, xtick distance = 1, ytick distance = 1, xmin=0, xmax=5, ymin=0, ymax=5, grid=major]
					\plot[blue, ->, semithick]coordinates {(0, 0) (2, 3)} node[right, fill=white, scale=0.75] {\(\begin{pmatrix}2 \\ 3\end{pmatrix}\)};
					\plot[purple, ->, semithick] coordinates {(0, 0) (1, 0)};
					\draw[black, ->, semithick] (axis cs: 1, 0) -- node[right, fill=white, scale=0.75] {\(\biggl\|\begin{pmatrix}2 \\ 3\end{pmatrix} - \begin{pmatrix}1 \\ 0\end{pmatrix}\biggr\| = \sqrt{10}\)} (axis cs: 2, 3);
				\end{axis}
			\end{tikzpicture}
		\end{figure}
	\end{enumerate}